% ADR/AD/Viewpoint technical handout
% Repo location: src/architecture/technical-documents/
% Document classification: technical reference; not a template/blank
% Standalone pdfLaTeX, no shell escape or external assets
\documentclass[10pt,letterpaper]{article}
\usepackage[T1]{fontenc}
\usepackage[utf8]{inputenc}
\usepackage{lmodern,microtype}
\usepackage[letterpaper,margin=.78in,headheight=16pt]{geometry}
\usepackage[table]{xcolor}
\usepackage{booktabs,tabularx,longtable,array,enumitem}
\usepackage[most]{tcolorbox}
\usepackage{fancyhdr,lastpage,titlesec}
\usepackage{hyperref,xurl}
\definecolor{Navy}{HTML}{193443}
\definecolor{Teal}{HTML}{137980}
\definecolor{Ink}{HTML}{293D47}
\definecolor{Muted}{HTML}{62737D}
\definecolor{Pale}{HTML}{F1F6F7}
\color{Ink}
\hypersetup{colorlinks=true,linkcolor=Teal,urlcolor=Teal}
\urlstyle{same}
\setlength{\parindent}{0pt}
\setlength{\parskip}{5pt}
\setlength{\emergencystretch}{2em}
\renewcommand{\arraystretch}{1.18}
\setlength{\tabcolsep}{4pt}
\setlist[itemize]{leftmargin=1.3em,topsep=3pt,itemsep=2pt}
\setlist[enumerate]{leftmargin=1.5em,topsep=3pt,itemsep=2pt}
\newcolumntype{Y}{>{\raggedright\arraybackslash}X}
\newcolumntype{L}[1]{>{\raggedright\arraybackslash}p{#1}}
\titleformat{\section}{\sffamily\Large\bfseries\color{Navy}}{\thesection}{.6em}{}
\titleformat{\subsection}{\sffamily\bfseries\color{Teal}}{\thesubsection}{.6em}{}
\pagestyle{fancy}
\fancyhf{}
\fancyhead[L]{\sffamily\scriptsize\color{Muted} ARCHITECTURE / DOCUMENT GOVERNANCE}
\fancyhead[R]{\sffamily\scriptsize\color{Muted} TECHNICAL HANDOUT}
\fancyfoot[L]{\sffamily\scriptsize\color{Muted} 29 October 2026 / v1.0}
\fancyfoot[R]{\sffamily\scriptsize\color{Muted} \thepage/\pageref*{LastPage}}
\renewcommand{\headrulewidth}{.3pt}
\newtcolorbox{principle}{breakable,colback=Pale,colframe=Teal,boxrule=.4pt,arc=1pt,left=8pt,right=8pt,top=6pt,bottom=6pt}
\begin{document}
{\sffamily\small\bfseries\color{Teal} REPOSITORY ARCHITECTURE DOCUMENTATION}\par
\vspace{7pt}
{\sffamily\fontsize{22}{25}\selectfont\bfseries\color{Navy}
Architecture Decisions, Descriptions, and Viewpoint Specifications\par}
\vspace{7pt}
{\large Standards-grounded classification, lifecycle, traceability, and audit}\par
\vspace{9pt}{\color{Teal}\hrule height 1.1pt}\vspace{8pt}
\begin{principle}
\textbf{The three-artifact rule.} An Architecture Decision Record explains
\emph{why a consequential choice was made}. An Architecture Description expresses
\emph{the architecture of an identified entity of interest for stakeholders and
their concerns}. An Architecture Viewpoint Specification defines
\emph{reusable conventions for creating, interpreting, and using a kind of
architecture view}. The three artifacts are related but not interchangeable.
\end{principle}

\section{Fact-checked terminology}
\textbf{Architecture versus its description.} ISO/IEC/IEEE 42010:2022
specifies requirements for the structure and expression of an
\emph{architecture description} (AD), not for the architecture itself.
The entity of interest (EoI) may be software, a system, an enterprise,
a product line, or another defined entity. The standard does not mandate a
specific authoring method, notation, medium, or tool.\textsuperscript{[1,2]}

\textbf{AD is the standards term.} The phrase
\emph{Architecture Description Record} is acceptable as a local repository
alias when explicitly defined; it is not an independently standardized
artifact type. An AD need not be a giant document or only an as-built model:
a versioned set of linked artifacts may express proposed, reference, or
realized architecture, with scope and applicability identified.

\textbf{Viewpoint is not view.} A \emph{viewpoint} specifies conventions
for framing concerns and constructing/interpreting views; a \emph{view}
is an information part of a particular AD. The 2022 standard specifies
separate requirements for architecture viewpoints and model kinds.
AVS is a convenient repository abbreviation, not mandatory ISO notation.

\textbf{ADR is a documented practice.} ADRs capture the context,
alternatives, choice, trade-offs, rationale, and consequences for one
architecturally significant decision. Nygard-style and MADR-style records
are practice conventions; the 42010 standard requires attention to
essential decisions and rationale within the AD, but does not mandate
independent ADR Markdown files.\textsuperscript{[1,3,4,5]}

\section{Key distinctions}
\small
\begin{tabularx}{\linewidth}{@{}L{.87in}Y Y Y@{}}
\toprule
& \textbf{ADR} & \textbf{AD} & \textbf{AVS}\\
\midrule
Question & Why this option? & What architecture is expressed
for this EoI and concern set? & How are views of this kind
created and interpreted?\\
Scope & One bounded decision & EoI, stakeholders, concerns,
viewpoints, views, correspondences, essential rationale &
Reusable concerns, model kinds, conventions, methods and interpretation\\
Form & Usually concise; arbitrary repository medium &
Versioned document or connected AD information parts &
Reusable viewpoint specification, often with examples\\
Change & Preserve accepted rationale; supersede material changes &
Update and baseline the description & Version conventions and
assess adopting views\\
Authority & Decision approver & AD owner/steward &
Viewpoint specification steward\\
Basis & Nygard/MADR/organizational policy &
ISO/IEC/IEEE 42010:2022 & ISO/IEC/IEEE 42010:2022\\
\bottomrule
\end{tabularx}
\normalsize
\begin{principle}
\textbf{Important limitations.} An ADR need not be exactly one page or
Markdown. An AD need not be long. An architectural diagram is not by itself
a conforming AD. Templates and PDFs never establish approval or
standards conformance automatically.
\end{principle}

\section{Lifecycle and change policy}
\begin{enumerate}
\item \textbf{ADR:} identify proposal, acceptance, rejection, deprecation
and supersession according to local policy. Accepted rationale should be
historically stable; create and link a superseding ADR when the decision
materially changes. Status updates and factual errata may be recorded
with an auditable history. This immutability convention is not an
ISO 42010 requirement.\textsuperscript{[3,4,5]}
\item \textbf{AD:} maintain EoI scope and version, identified concerns
and stakeholders, governing viewpoints, views, model kinds,
correspondences, known inconsistencies, decisions, rationale and
evidence. Explicitly distinguish current, proposed and reference baselines.
\item \textbf{AVS:} version the viewpoint identity, owner, scope,
concerns framed, typical stakeholders, model kinds, construction
methods, interpretation rules and correspondence methods. Review
dependent views whenever conventions change.
\end{enumerate}

\section{Traceability example}
\begin{principle}
\centering\sffamily\small
\textbf{AD-PLATFORM-001} (system and concern baseline)\\
$\downarrow$ uses \quad \textbf{VP-DEPLOYMENT-001} (viewpoint conventions)\\
$\downarrow$ governs \quad \textbf{VIEW-DEPLOYMENT-001} (system-specific view)\\
$\leftrightarrow$ justified by \quad \textbf{ADR-0042} (choice and trade-offs)\\
$\downarrow$ verified by \quad \textbf{REQ-18 / TEST-27 / EVIDENCE-09}
\end{principle}
For a cloud deployment, a resilience concern in the AD motivates a deployment
view. An AVS establishes representation and interpretation conventions. An ADR
records why multi-zone deployment was chosen over competing options. Evidence
establishes whether the accepted choice was implemented. A hyperlink
provides traceability, not proof that two models are consistent.

\section{Findings: unresolved and hybrid repository artifacts}
The existing artifact registry already distinguishes ten ADR
Template--Blank pairs, one AD pair, and a legacy viewpoint pair. The
following review of additional sources does not assert that new Blanks
have already been created.
\small
\begin{longtable}{@{}L{1.05in}L{1.22in}L{3.33in}@{}}
\toprule
\textbf{Source}&\textbf{Classification}&\textbf{Disposition / next action}\\
\midrule\endhead
Polished 42010 VP &
Viewpoint authoring guide and worked example &
Preserve the instructional guide, specification outline and
service-continuity illustration. Avoid generating a pointless
blank of explanatory material; extract any uniquely reusable form
only when needed.\\
\addlinespace
Polished VP worksheet &
AVS Template, currently unpaired &
Create a printable companion Blank preserving identity, stakeholder
concerns, model-kind catalogue, view methods, correspondence
methods, review fields, and attribution. The accompanying
viewpoint section is a non-standalone source fragment.\\
\addlinespace
42010 document 02 &
Composite AD register suite &
Keep AD identity, stakeholders, concerns, perspectives, aspects,
viewpoints, views, view components, correspondences, decision/rationale
and review registers coherent. Pair the suite or extract independently
reusable registers; do not classify the entire suite as an ADR.\\
\addlinespace
Cloud forms T01--T08 &
Mixed architecture instruments &
T01 context (AD input); T02 deployment-model decision (decision
worksheet); T03 quality scenarios; T04 mechanisms; T05 ADR worksheet;
T06 recovery; T07 architecture review; T08 traceability. Pair each
extracted instrument according to its actual purpose.\\
\addlinespace
Security forms S01--S08 &
Security governance instruments &
S01 threat boundaries, S02 IAM, S03 network, S04 data,
S05 hardening, S06 logging, S07 verification/exception and
S08 incident readiness. These may inform an AD but are not
automatically ADs, ADRs, or AVSs.\\
\bottomrule
\end{longtable}
\normalsize
\textbf{Full reviewed paths:}
\path{src/architecture/templates/42010-vp-template-polished.tex};
\path{src/architecture/templates/viewpoint-template-polished.tex};
\path{src/architecture/standards/iso-iec-ieee-42010/02-ad-templates.tex};
\path{src/architecture/cloud/engineering-suite/05-architecture-documentation-templates.tex};
\path{src/architecture/cloud/engineering-suite/06-security-governance-templates.tex}.

\section{Template--Blank lifecycle and validators}
A Template is an instructional authoring instrument; a Blank is a distinct,
printable, uncompleted working copy. A completed example, technical guide,
answer key, or included source fragment is not automatically a Template.
\begin{enumerate}
\item \textbf{Common:} unique pair ID, authoritative paths, equivalent required
structure, writing fields, PDF build, leaf README, and no visible square-bracketed
placeholder words in the Blank. Preserve TeX package options and optional
arguments that use square brackets syntactically.
\item \textbf{ADR-specific:} singular choice, status, alternatives, rationale,
consequences, accountable authority, revisit/supersession and AD/evidence links.
\item \textbf{AD-specific:} entity and baseline, stakeholder concerns,
viewpoints and views, components/model kinds, correspondences,
inconsistencies, essential decisions and rationale.
\item \textbf{AVS-specific:} viewpoint identity/version, intended concerns
and typical stakeholders, conventions and model kinds, view/interpretation
methods, correspondence methods, steward, sources and limitations.
\item \textbf{Composite:} retain a coherent register suite where the
sections are interdependent; otherwise extract independent forms and
give each an ID and Template--Blank pair. Preserve attribution and
original worked examples. Do not fabricate a completed record.
\end{enumerate}

\section{Integration into the LaTeX repository}
\begin{tabularx}{\linewidth}{@{}L{1.03in}Y@{}}
\toprule
\textbf{Location}&\textbf{Contract}\\
\midrule
Technical documents & This handout is a technical reference, not a pair member.\\
ADR templates & Decision Template--Blank pairs and governance lifecycle.\\
Architecture templates & AD form, AVS worksheet, VP guide and source
fragment remain visibly distinct.\\
ISO 42010 suite & Preserve standards reference, examples and composite AD
registers.\\
Cloud engineering & Classify T01--T08 and S01--S08 by each form's purpose.\\
Pair registry & Only completed pairs go under the pairs array;
unresolved sources remain in the review queue with explicit dispositions.\\
Validator & Test each family semantically and check rendered PDFs;
a structural pass is not conformance.\\
\bottomrule
\end{tabularx}
\textbf{Repository checks:} \texttt{make preflight};
\texttt{make build-category-architecture};
\texttt{python3} \path{tooling/scripts/validate_artifact_pairs.py}.

\section{Sources and normative caution}
\small
\begin{enumerate}[label={[\arabic*]},leftmargin=2em]
\item ISO, \emph{ISO/IEC/IEEE 42010:2022}, second edition.
\url{https://www.iso.org/standard/74393.html}.
\item IEEE Standards Association, \emph{IEEE/ISO/IEC 42010-2022}.
\url{https://standards.ieee.org/ieee/42010/6846/}.
\item Michael Nygard, \emph{Decision record template}.
\url{https://github.com/architecture-decision-record/architecture-decision-record/blob/main/architecture-decision-record.github.io/src/content/templates/decision-record-template-by-michael-nygard.md}.
\item AWS Prescriptive Guidance, \emph{Architectural decision record process}.
\url{https://docs.aws.amazon.com/prescriptive-guidance/latest/architectural-decision-records/adr-process.html}.
\item Microsoft Learn, \emph{Maintain an architecture decision record}.
\url{https://learn.microsoft.com/en-us/azure/well-architected/architect-role/architecture-decision-record}.
\end{enumerate}
\textbf{Authority and limits.} The ISO/IEEE official abstracts support
the standard's published scope. Clause-level conformance assessment
requires the controlling normative text. Local ADR lifecycle, naming,
versioning, and printable Template--Blank governance are
repository rules, not additional ISO requirements.
\end{document}
